\documentclass[11pt]{article}
\usepackage[margin=0.7in]{geometry}
\usepackage{xcolor}
\usepackage{enumitem}
\usepackage{hyperref}
\setlength{\parindent}{0pt}
\setlength{\parskip}{6pt}
\sloppy
\newcommand{\loc}[1]{\textbf{Location:} #1}
\newcommand{\current}[1]{\textbf{Verified evidence:} #1}
\newcommand{\replacewith}[1]{\textbf{Replace with / Required action:}\\{\color{blue}#1}}
\newcommand{\why}[1]{\textbf{Audit conclusion:} #1}
\begin{document}
\begin{center}
{\Large Verified Analysis Audit and Mandatory Revision Guide}\\
{\small Data leakage, ablation, and statistical testing audit. Blue text indicates required replacement or action.}
\end{center}

\section*{Audit Conclusion}
The current performance claims cannot yet be treated as independently validated. The documented workflow contains a direct leakage contradiction, no component-cost ablation is reported, and the Wilcoxon p-values show that correlated results were almost certainly pooled beyond the 10 independent random seeds. The designated GitHub repository has one branch and four commits containing only submission and revision materials; no analysis code, dataset, seed-level results, or embedded PDF attachments were available for rerunning the models. The findings below therefore distinguish (i) defects proven from the manuscript and reported numbers, (ii) analyses that must be rerun, and (iii) text that may be used only after the rerun passes.

\section*{Edit 1. Remove the unsupported no-leakage assertion}
\loc{Page 7, Figure 1; page 10, Table 3; page 11, lines 206--209}
\current{Figure 1 explicitly orders cut-off analysis and Min-Max scaling before the 70:30 train-test split. Table 3 reports thresholds and class counts for all 419 projects: mean = USD 148,271, Q3 = USD 180,676, and Q3 + IQR/2 = USD 228,626. The mean threshold is the rounded full-dataset target mean in Table 2. Page 11 nevertheless states that preprocessing and model selection used training data only. These records are mutually inconsistent.}
\replacewith{Delete the present lines 206--209 until the code has been checked or the analysis rerun. In Figure 1, move ``Train-test split (70:30; 10 random seeds)'' ahead of cut-off calculation, encoding, scaling, feature selection, and GridSearchCV. For every outer split, calculate target-derived thresholds from the outer training partition only. Fit the encoder and scaler inside a Pipeline within the training/inner-validation process, never on the full dataset. Apply the fitted preprocessing objects and training-derived threshold unchanged to the outer test partition.}
\why{As drawn and tabulated, the test targets influence the cost threshold and the test features influence scaling before evaluation. This is leakage unless Figure 1 and Table 3 are merely descriptive and the unavailable code proves a different order. The current sentence cannot be retained as evidence of a procedure that the workflow contradicts.}

\section*{Edit 2. Restore a complete split and model-selection protocol}
\loc{Page 11, immediately before Section 3.2}
\current{The revised manuscript no longer gives the full split protocol in prose. Figure 1 is the only place showing a 70:30 split and 10 seeds. GridSearchCV is named, but the number of folds, scoring rule, pipeline order, stratification, seed values, and validation role are absent.}
\replacewith{For each of 10 prespecified random seeds, the raw dataset was divided into an outer training partition (70\%) and an outer test partition (30\%). All target-derived cutoffs, encoders, scalers, and model hyperparameters were estimated without access to the outer test partition. Hyperparameter tuning was performed by [K]-fold cross-validation within the outer training partition using [primary scoring metric]. Preprocessing was included inside each cross-validation fold. After selection, the model was refitted on the complete outer training partition and evaluated once on the corresponding outer test partition. Report the 10 seed values, stratification rule, K, scoring metric, and complete search spaces.}
\why{Fill every bracket from the actual rerun. Without these details, the leakage-control claim is neither auditable nor reproducible.}

\section*{Edit 3. Prevent test-set model-selection leakage}
\loc{Page 7, lines 161--164; pages 18--25; SHAP model selection}
\current{The workflow says that the ``best-performing model'' was selected and then interpreted. Nine cutoff-by-transfer configurations are compared on the repeated outer test sets, and the mean-cutoff balanced model is chosen from those test results for SHAP. This reuses the evaluation data for configuration selection.}
\replacewith{Choose the final cutoff and transfer strategy using training-only inner validation within each outer split, or pre-specify one primary configuration before test evaluation. If all nine configurations are retained as an exploratory comparison, do not call the test-selected configuration independently validated and do not use its test advantage as confirmatory evidence. Apply SHAP to a prespecified model or explain that the interpretation is exploratory.}
\why{Holding out a test set is insufficient if the same test results are then used to choose the preferred framework. A nested or prespecified selection rule is required for an unbiased final-performance claim.}

\section*{Edit 4. Replace Table 3 with split-specific evidence}
\loc{Page 10, Table 3; page 22, Figure 9}
\current{The threshold values and 243/176, 314/105, and 352/67 class counts use the complete sample of 419 projects. They do not show what each training-only cutoff or outer test composition was across the 10 seeds.}
\replacewith{Retitle the existing table ``Full-dataset descriptive thresholds; not used for model fitting'' if it is kept. Add a reproducibility table reporting, across the 10 outer splits, the mean, standard deviation, minimum, and maximum of each training-derived threshold; training low/high counts; test low/high counts; and the number of high-cost observations used for each segment-wise RMSE.}
\why{A fixed full-sample target threshold can leak test-target information. Split-specific reporting also reveals whether high-cost RMSE is based on too few observations in some test splits.}

\section*{Edit 5. Run the component-cost ablation}
\loc{Pages 8--9, Table 2; Abstract; Implications; Conclusions}
\current{N10, N11, N14, and N15 are monetary component costs for windows, doors, HVAC, and boilers, while the output is total direct construction cost. Their reported mean sum is USD 44,851.88, approximately 30.3\% of the target mean, and N10 alone has r = 0.81 with the target. No ablation or provenance/timing analysis is reported.}
\replacewith{Rerun the complete nested evaluation on the same outer splits for three feature sets: (A) all 21 inputs; (B) all inputs except N10, N11, N14, and N15; and (C) physical, categorical, and quantity variables only. For every set, report overall R2, MAPE, RMSE, high-cost RMSE, and paired seed-level differences with confidence intervals. State whether each monetary input is a preliminary estimate available at the claimed decision point or a realized component of the final direct construction cost.}
\why{If the component costs are contained in the target or are known only after detailed costing, the task is partly total-cost reconstruction rather than early-stage estimation. The present data are insufficient to quantify the effect; the ablation is mandatory.}

\section*{Edit 6. Narrow the claim if the ablation cannot be run}
\loc{Title, Abstract, Introduction, Implications, and Conclusions}
\current{The manuscript repeatedly claims reliable early-stage cost estimation without defining when the component-cost inputs become available.}
\replacewith{If N10, N11, N14, and N15 cannot be shown to be available at the early planning stage, replace ``early-stage cost estimation'' with ``conditional retrofit cost estimation using partial component-cost information.'' Add this limitation: ``Because several predictors are component-level cost estimates, the reported performance may not transfer to settings where only physical and quantity information is available.''}
\why{The stronger early-stage claim is supportable only after input timing is documented and the no-component-cost ablation remains practically useful.}

\section*{Edit 7. Withdraw Tables 5--6 in their current inferential form}
\loc{Pages 20 and 22, Tables 5--6}
\current{The paper states that performance was averaged over 10 random seeds. For 10 nonzero pairs, the smallest possible two-sided exact Wilcoxon p-value is 2/2^10 = 0.001953. Table 5 reports 1.86E-09, which equals the theoretical minimum 2/2^30 to rounding. This is compelling numerical evidence that 30 pairs were used, most plausibly 10 seeds multiplied by three transfer strategies. Table 6 similarly groups results across three cutoffs. Multiple observations from the same seed are correlated and are not 30 independent experimental replications.}
\replacewith{Do not retain the current p-values as confirmatory evidence. Reanalyze the 3-cutoff by 3-transfer-strategy design with seed as the independent block. Preferred analysis: for each metric, fit a repeated-measures or mixed-effects model with fixed effects for cutoff, transfer strategy, and their interaction, and a random intercept for seed; use a seed-clustered permutation or bootstrap if distributional assumptions are doubtful. Report the interaction before pairwise contrasts. Adjust planned contrasts with Holm's method and report paired effect estimates and 95\% confidence intervals, not p-values alone.}
\why{Pooling three correlated model results per seed produces pseudoreplication, artificially increases the nominal sample size from 10 to 30, and can make p-values much smaller than the independent-seed design supports.}

\section*{Edit 8. Correct the significance narrative}
\loc{Page 22, lines 453--461; Table 6; Conclusions}
\current{The manuscript claims that balanced transfer significantly improves both high-cost RMSE measures. Even before correcting the n=30 dependence problem, Holm adjustment across the 15 Table 6 tests changes conservative-versus-balanced RMSE above Q3 from p = 0.0293 to 0.1758. The Q3 + IQR/2 contrast changes from p = 0.00322 to 0.03289 under the same naive adjustment. Overall R2 and RMSE for conservative versus balanced are already nonsignificant at p = 0.158 and 0.146.}
\replacewith{Until the seed-blocked reanalysis is complete, replace the inferential claim with: ``Descriptively, balanced transfer produced lower average RMSE in the two high-cost subsets, whereas conservative transfer produced marginally better overall R2, MAPE, and RMSE. The present analysis does not establish universal statistical superiority of either strategy.''}
\why{The current broad claim is unsupported by the reported overall tests, uncorrected multiple comparisons, and the pooled analysis.}

\section*{Edit 9. Correct the best-model statement}
\loc{Abstract; pages 18--19; Table 4; Implications; Conclusions}
\current{Table 4 shows that mean-cutoff conservative transfer has the best overall values: R2 = 0.9568, MAPE = 10.79\%, and RMSE = USD 21,360. Mean-cutoff balanced transfer has R2 = 0.9566, MAPE = 10.95\%, and RMSE = USD 21,377, but lower high-cost RMSE. The overall RMSE difference is only USD 17.}
\replacewith{Describe conservative transfer as the descriptive overall-metric leader and balanced transfer as the descriptive high-cost-RMSE leader. Do not select balanced as the single best model unless a primary high-cost endpoint and selection rule were specified without using the outer test results. Update every occurrence of R2 = 0.9566 and the 46\%/25\% reductions after the leak-free rerun.}
\why{The current narrative chooses one configuration after inspecting different endpoints and then reports it as the best-performing framework.}

\section*{Edit 10. Add code, data, and seed-level result availability}
\loc{End matter, before References}
\current{No Data Availability or Code Availability statement is present. The accessible submission repository contains no analysis scripts, dataset, fitted preprocessing objects, seed-level metrics, or ablation outputs, and the manuscript PDF has no embedded attachments.}
\replacewith{Data and code availability: Deidentified input data or a governed-access procedure, complete preprocessing and modeling code, package versions, the 10 random seeds, all hyperparameter grids, split indices, training-derived thresholds, and a tidy seed-level metrics file are available at [permanent repository and DOI]. The seed-level file contains one row per seed, cutoff, transfer strategy, feature-set ablation, and metric.}
\why{This is necessary to verify the corrected leakage controls, ablation, and repeated-measures analysis. Do not state public availability until the archive actually exists.}

\section*{Minimum Seed-Level Results Schema}
\loc{Required supplementary CSV}
\current{Table 4 reports means and standard deviations for only some metrics, while Tables 5--6 depend on unreported paired observations.}
\replacewith{Include these columns: seed; split ID; train/test indices or hashes; cutoff rule; training-derived cutoff value; transfer strategy; feature set; train/test low/high counts; classifier precision/recall/F1; R2; MAPE; overall RMSE; RMSE at or above Q3; RMSE at or above Q3 + IQR/2; selected hyperparameters; frozen layers; learning rate; and run status.}
\why{Without the paired seed-level file, no reviewer can reproduce the p-values, confidence intervals, or ablation comparisons.}

\section*{Submission Gate}
\loc{Before resubmission}
\current{The current PDF contains a leak-free assertion that conflicts with its workflow and confirmatory claims based on an invalidly pooled statistical analysis.}
\replacewith{Do not submit the current performance and significance claims unchanged. First rerun the pipeline with split-first preprocessing and training-only thresholds; use training-only configuration selection; complete the component-cost ablation; analyze seed-blocked repeated measures with multiplicity control; then update the abstract, Table 3, Tables 4--6, SHAP model choice, implications, conclusion, and cover letter from the corrected outputs.}
\why{Text-only edits cannot repair leakage, missing ablation, or pseudoreplication. The current numerical results should be labeled provisional until the corrected run is complete.}

\end{document}
